\documentclass[11pt,a4paper]{article}
\usepackage[margin=25mm]{geometry}
\usepackage{amsmath,amssymb,amsthm,booktabs,hyperref}
\setlength{\emergencystretch}{2em}
\newtheorem{theorem}{Theorem}
\newtheorem{lemma}{Lemma}
\title{Identifying the historical quartic input set}
\author{Human Compression project: a component of the L059 recovery}
\date{}
\begin{document}\maketitle

The numerical work of De Saedeleer, Leemans, Mixer and Pisanski
\cite{l059} uses the census of edge-transitive tetravalent graphs
of Wilson and Poto\v cnik. The census has since been expanded and its
identifiers have changed. Here we identify the set of graph isomorphism
classes required for a corrected version of that calculation. We use
an explicit coverage assertion of the census authors, together with
finite checks of the retained data. We do not reconstruct the meaning
of every historical identifier.

All graphs are finite, simple and connected. A tetravalent graph has
degree four at each vertex. For a connected bipartite graph $X$, let
$F=\operatorname{Aut}(X)$ and let $A=\operatorname{Aut}^{+}(X)$ be
the subgroup preserving each part of its bipartition. A graph is
\emph{dart-transitive} if $F$ is transitive on ordered adjacent pairs;
it is \emph{semisymmetric} if $F$ is edge-transitive but not
vertex-transitive. Define the property $P$ by
\[
 |V(X)|\leq150,\qquad X\text{ bipartite},\qquad
 A\text{ edge-transitive},\qquad |F_v|<200.
 \tag{1}
\]
In this setting the order of $F_v$ is independent of the choice of
$v$. In the semisymmetric case the two vertex orbits are the two
parts, which have equal size because the graph is regular.

\begin{lemma}
Among bipartite edge-transitive graphs, $A$ is edge-transitive
exactly in the dart-transitive and semisymmetric cases. In these
cases
\[
                    \frac{|A|}{|V(X)|}=\frac{|F_v|}{2}.
\]
\end{lemma}
\begin{proof}
If $F$ is not vertex-transitive, its two vertex orbits are the two
parts. Thus $A=F$ is edge-transitive and orbit--stabilizer gives
$|F|=(|V(X)|/2)|F_v|$. If $F$ is dart-transitive, any edge may be
sent to another with its black endpoint sent to the black endpoint.
Such an automorphism belongs to $A$. Here $[F:A]=2$, and
$|F|=|V(X)|\,|F_v|$.

In the remaining case $F$ is vertex- and edge-transitive but not
dart-transitive. If $A$ were edge-transitive, its point stabilizer
would be transitive on the four neighbours of a point. Every
automorphism fixing a point preserves the parts, so $A_v=F_v$.
Vertex-transitivity of $F$ would then imply dart-transitivity,
a contradiction. The two stated formulas follow from the two
orbit--stabilizer calculations.
\end{proof}

\paragraph{The acquired premises.}
We retain two historical copies of the census table, archived on
5 May 2008 and 28 December 2014. The historical homepage describes
the table as the list of graphs in the census. As with any use of
this census, we acquire the authors' assertion that its entries
represent distinct isomorphism classes and that the historical
graph invariants are correct. We do not claim an independent
reconstruction of every historical graph.

The later primary source \cite[\S1, physical p.~2]{recipes} describes
the expanded census as
\begin{quote}
``a collection of all known edge-transitive graphs of valence 4
up to 512 vertices.''
\end{quote}
We acquire precisely the resulting coverage premise: every graph
in either retained historical census occurs, up to isomorphism,
in the later census at the linked site. These historical graphs
were already known before the 2016 statement. The retained current
homepage cites that paper, links the full Magma data, and explains
that the older identifiers changed when the search was reorganized.
The files, URLs and hashes are recorded in \texttt{inputs.json} and
\texttt{source\_acquisition.json}.

This is a coverage premise about graphs already in an earlier
catalogue. The authors explicitly do not assert completeness among
all semisymmetric tetravalent graphs. We acquire no such stronger
classification.

\paragraph{Finite verification.}
The checker parses every row of each of the three tables and every
graph occurrence in the full Magma source. The results are:
\begin{center}
\begin{tabular}{lrrrr}
\toprule
Edition & Table rows & Rows satisfying $P$ & Dart-transitive & Semisymmetric\\
\midrule
5 May 2008 & 788 & 397 & 272 & 125\\
28 December 2014 & 793 & 397 & 272 & 125\\
Current full census & 7364 & 397 & 272 & 125\\
\bottomrule
\end{tabular}
\end{center}
There are 807 graphs of order at most 150 in the current data,
of which 562 are bipartite. All 562 are checked to be connected,
simple, tetravalent and pairwise nonisomorphic, using canonical
nauty certificates. Their full automorphism groups, edge orbits,
vertex orbits and, when needed, dart orbits are computed. Each
computed automorphism is also checked directly against the edge
set. The resulting symmetry type, full group order and vertex
stabilizer order agree with the current table in every case.
Exactly 397 satisfy (1). Thus the current graph count is checked
on actual edge sets, independently of the printed count.

\begin{theorem}
Under the stated historical census and coverage premises, either
retained historical census and the retained current census contain
the same 397 graph isomorphism classes satisfying $P$.
\end{theorem}
\begin{proof}
Let $\mathcal H$ be the set of classes satisfying $P$ in one of
the historical editions, and let $\mathcal C$ be the corresponding
set in the current census. The acquired historical assertion and
the parsed row count give $|\mathcal H|=397$. The computations on
current edge sets give $|\mathcal C|=397$.

Every graph representing a class in $\mathcal H$ was known before
the expanded census and has at most 150 vertices, hence at most
512. The coverage premise places an isomorphic copy in the later
census. Property $P$ is preserved by graph isomorphism: an
isomorphism conjugates the full automorphism groups, carries the
two bipartition classes to the two classes, and preserves the
orders and orbit conditions in (1). Thus
$\mathcal H\subseteq\mathcal C$. Since the sets have the same
finite cardinality, $\mathcal H=\mathcal C$. The argument applies
separately to the two historical editions.
\end{proof}

\paragraph{Scope and attribution.}
The threshold in (1) is explicit. The original paper says that the
vertex stabilizer has order less than 100. That is not the present
threshold: by the lemma, (1) is instead equivalent to
$|\operatorname{Aut}^{+}(X)|/|V(X)|<100$. The original paper also
explicitly includes $K_{4,4}$, whose full vertex stabilizer has
order $144$. Thus its stated cutoff and its stated example already
disagree. The present theorem identifies the historical input set
for the corrected cutoff; it neither silently changes the
original statement nor establishes why the discrepancy arose.

The graph identification theorem concerns uncoloured isomorphism
classes. It does not fix an ordering of the two parts, transfer a
particular historical identifier, or certify the later marking
enumeration. Those are separate tasks. The census construction,
historical invariants and coverage assertion are borrowed from
Wilson and Poto\v cnik. Our contribution here is the explicit
finite comparison and the short set argument permitting reuse
of the current edge data. No new classification or solution of
an open problem is claimed.

\begin{thebibliography}{9}
\bibitem{l059} J.~De Saedeleer, D.~Leemans, M.~Mixer and T.~Pisanski,
\emph{Core-free, rank two coset geometries from edge-transitive
bipartite graphs}, Mathematica Slovaca \textbf{64} (2014), 991--1006;
arXiv:1106.5704v2.
\bibitem{recipes} S.~Wilson and P.~Poto\v cnik,
\emph{Recipes for Edge-Transitive Tetravalent Graphs},
arXiv:1608.04158v1 (2016), \S1; accompanying
\href{https://jan.ucc.nau.edu/swilson/C4FullSite/index.html}{full census}.
\bibitem{nauty} B.~D.~McKay and A.~Piperno,
\emph{Practical graph isomorphism, II}, Journal of Symbolic
Computation \textbf{60} (2014), 94--112. The finite checker uses
nauty through the pynauty interface; its installed version is
recorded in the receipt.
\end{thebibliography}
\end{document}
